% =============================================================
% capitoli/capXX_argomento.tex
% Template per un nuovo capitolo
%
% Istruzioni:
%   1. Rinominare il file (es. cap01_numeri_naturali.tex)
%   2. Sostituire tutti i segnaposto [...]
%   3. Includere in main.tex con \input{capitoli/capXX_argomento}
% =============================================================

\chapter{[Titolo del Capitolo]}
\label{cap:[etichetta]}

% -------------------------------------------------------------
% OBIETTIVI DI APPRENDIMENTO
% -------------------------------------------------------------
\begin{tcolorbox}[
    colback  = gray!6,
    colframe = colorePrimario,
    arc      = 4pt,
    title    = {\bfseries\color{colorePrimario} Obiettivi di apprendimento},
    fonttitle = \bfseries
]
Al termine di questo capitolo sarai in grado di:
\begin{itemize}
    \item [Obiettivo 1]
    \item [Obiettivo 2]
    \item [Obiettivo 3]
\end{itemize}
\end{tcolorbox}

\spazio

% -------------------------------------------------------------
% SEZIONE 1 — Prima sezione del capitolo
% -------------------------------------------------------------
\section{[Titolo sezione]}
\label{sec:[etichetta]}

[Testo introduttivo della sezione. Esporre il concetto con linguaggio
chiaro e diretto. Definire ogni termine tecnico prima di utilizzarlo.
Frasi brevi. Nessuna ambiguità.]

% --- Definizione ---
\begin{definizionesemplice}{[Nome della definizione]}
    [Testo della definizione. Preciso, formale, completo.]
    \[
        [Formula\ se\ necessaria]
    \]
\end{definizionesemplice}

\spazio

% --- Regola / Proprietà ---
\begin{regola}{[Nome della regola o proprietà]}
    [Enunciato della regola. Eventuale formula:]
    \[
        [Formula]
    \]
\end{regola}

\spazio

% --- Esempio risolto ---
\begin{esempio}{[Titolo breve dell'esempio]}
    [Testo del problema / situazione di partenza.]

    \medskip
    \textbf{Soluzione:}

    \begin{align*}
        [Passaggio\ 1] &= [risultato] && \text{[spiegazione passaggio 1]} \\
        [Passaggio\ 2] &= [risultato] && \text{[spiegazione passaggio 2]} \\
        [Risultato\ finale] &= [valore]
    \end{align*}

    [Commento finale sul risultato, se necessario.]
\end{esempio}

\spazio

% --- Errore comune ---
\begin{attenzione}[Errore comune]
    [Descrizione dell'errore tipico che gli allievi commettono.
    Mostrare il ragionamento errato e quello corretto.]

    \medskip
    {\color{red}\textbf{Sbagliato:}} $[espressione\ errata]$

    \medskip
    {\color{colorePrimario}\textbf{Corretto:}} $[espressione\ corretta]$
\end{attenzione}

\spazio

% --- Lo sapevi? ---
\begin{sapevi}
    [Curiosità storica, applicazione reale o collegamento
    interdisciplinare legato all'argomento del capitolo.]
\end{sapevi}

\spazio

% -------------------------------------------------------------
% SEZIONE 2 — (aggiungere sezioni secondo necessità)
% -------------------------------------------------------------
% \section{[Titolo sezione 2]}

% -------------------------------------------------------------
% RIEPILOGO DEL CAPITOLO
% -------------------------------------------------------------
\section*{Riepilogo}

\begin{riepilogo}
    \begin{itemize}
        \item [Concetto chiave 1: definizione sintetica]
        \item [Concetto chiave 2: definizione sintetica]
        \item [Regola principale: enunciato sintetico]
        \item [Eventuali formule da ricordare]
    \end{itemize}
\end{riepilogo}

% =============================================================
% Fine capitolo
% =============================================================
